\chapter{The Simulation Program}
\label{chap:sims}

A framework earns the name only if it can be computed and confronted, and this
one's strength is that it can: rare vacuum fluctuations, Einstein--Cartan bounces,
and thermodynamic equilibrium are all tractable problems. The path from
``interesting picture'' to ``constrained by data'' runs through simulation, and it
is tiered---each tier a project of definite scope, higher tiers building on lower.
This chapter lays out that pathway and, more importantly, reports what already
runs: the bounce, the adiabatic extrusion that protects the inherited matter
excess, and a scale-invariant primordial spectrum with the right acoustic peaks,
all reproducible from the code in \texttt{sims/}.

\section{Phase 0: is our universe a black-hole interior?}
\label{sec:phase0}

Before any dynamical run there is cheaper, sharper work to do. The thesis under
everything is that our universe is the interior of a (rotating) black hole behind
a bounce membrane; in the Copernican spirit we do not assume it but ask what the
universe's \emph{measured} numbers already imply. Four ``corners,'' needing only
arithmetic and existing catalogues, each tie an observed quantity to a black-hole
property (\Cref{fig:firstlight}).

\paragraph{A --- flatness is the Schwarzschild condition.} For a Friedmann
universe the Schwarzschild radius of the mass inside the Hubble sphere works out
to $R_s = \Omega\,R_H$ \emph{identically}, so the measured flatness
$\Omega = 1.000\pm0.005$ \emph{is} the statement that the observable universe sits
at its own Schwarzschild radius---with Planck values
$M\approx9\times10^{52}\,\mathrm{kg}$ and $R_s/R_H=1.000$. A necessary
consistency condition, not a proof, but a striking one.

\paragraph{B --- the spin is sub-extremal, and predicts an axis.} The largest
causal angular momentum of that mass gives a dimensionless spin
$a^*_{\max}=2/\Omega\approx2$, so a realistic small net rotation leaves the
interior comfortably sub-extremal---a censored Kerr interior, not a naked
singularity. The observable handle is \emph{direction}: a parent-inherited axis
should be shared by the CMB ``axis of evil'' and the galaxy-spin
asymmetry~\citep{shamir2025}, the cheapest potentially-decisive test in the whole
program (Chapter~\ref{chap:nowobs}).

\paragraph{C --- the background cannot be the \emph{present} horizon's glow.} The
Hawking temperature of the Hubble-sphere mass is $\sim\!10^{-30}\,\mathrm{K}$,
thirty orders below the $2.7\,\mathrm{K}$ we measure. So the microwave background
can only be the \emph{parent's} horizon radiation, hot at emission near the bounce
and redshifted by our expansion---which means a successful calculation must
deliver a spectrum, not merely a temperature (Tier~2 below).

\paragraph{D --- the dark sector fixes one filter number.} If dark matter is the
membrane-rejected counterpart matter, the dark-to-baryon ratio fixes the
pass-fraction $f=\omega_b/(\omega_b+\omega_c)$; Planck's $\omega_c/\omega_b\approx
5.4$ gives $f\approx1/6$. Any Einstein--Cartan filter model must reproduce that one
number---the tightest quantitative target the membrane physics faces.

\begin{figure}[t]
  \centering
  \includegraphics[width=\textwidth]{bh_first_light}
  \caption{Phase-0 consistency checks, read straight off Planck cosmology.
    \emph{Left:} the Hawking temperature $T_H(M)$ puts the observable universe
    $\sim$30 orders below the CMB, so the CMB is not present-horizon radiation.
    \emph{Right:} the black-hole identities $R_s/R_H=\Omega\approx1$, maximal spin
    $a^*_{\max}\approx2$, and membrane filter $f\approx0.16$. Reproducible as
    \texttt{figures/scripts/ch02\_bh\_first\_light.py}.}
  \label{fig:firstlight}
\end{figure}

Phase~0 establishes that the \emph{kinematics} are consistent with a black-hole
interior. The tiers establish the \emph{mechanism}.

\section{Tier 1: the bounce itself}

The first dynamical task is to evolve a spherical collapse of fermionic matter in
Einstein--Cartan gravity and watch torsion halt it. It works: a homogeneous
collapse reaches the bounce density and reverses where general relativity would
hit a singularity (the result already shown as \Cref{fig:bounce}), and the bounce
is \emph{robust}---it occurs across the equation of state, not for a fine-tuned
fluid. Gluing the bouncing interior to a Schwarzschild exterior---the 1939
Oppenheimer--Snyder construction with the singularity replaced by a bounce---shows
both faces at once: a distant observer sees the surface freeze at the horizon and
redshift to black (a black hole), while inside the matter reaches the bounce
density at a tiny radius, bounces, and re-expands into a child universe. A custom
finite-difference code on a radial grid is enough; no heavy numerical-relativity
machinery is required to demonstrate the core phenomenon.

\section{Tier 1c: the extruder is adiabatic}
\label{sec:extruder}

The single most consequential numerical result so far concerns not the bounce but
its \emph{entropy budget}, because that is what protects the inherited matter
excess of Chapter~\ref{chap:dawn}. Entropy is generated only by irreversibility,
and the bounce is far too fast to thermalise: the controlling small parameter is
the same $\Omega_{\mathrm{bounce}}/T_{\mathrm{bounce}}\sim1.5\times10^{-11}$ that
limits the chiral-vortical estimate. Computing the dilution factor
$D=S_{\mathrm{after}}/S_{\mathrm{before}}$ from the dynamics, in both the
homogeneous bulk-viscous channel and the inhomogeneous shear (BKL/Mixmaster)
channel, gives $D=1$ to about one part in $10^{11}$ (\Cref{fig:extruder}). The
shear result is the more surprising: in plain general relativity anisotropy
diverges into the singularity and the entropy production runs away, but the same
torsion that removes the singularity \emph{caps the density}, bounces the shear,
and keeps the entropy integral finite. The conclusion is sharp:
$\eta_{\mathrm{child}}=\eta_{\mathrm{parent}}$, inheritance is essentially exact,
and the factor-of-$41$ shortfall of the chiral-vortical estimate binds only the
parentless first universe---every descendant simply carries down the excess made
once at the top.

\begin{figure}[t]
  \centering
  \includegraphics[width=\textwidth]{extruder}
  \caption{The bounce entropy budget. \emph{Left:} the entropy-production kernel
    peaks on the shoulders of the bounce and vanishes at the turnaround; the total
    is tiny. \emph{Right:} the dilution $D-1$ versus bulk viscosity stays
    $\lesssim10^{-10}$ across the whole physical range, so
    $\eta_{\mathrm{child}}\approx\eta_{\mathrm{parent}}$. Reproducible as
    \texttt{sims/src/cartasis\_sims/extruder.py},
    \texttt{figures/scripts/ch09\_extruder.py} (and \texttt{ch09\_shear.py} for the
    shear channel).}
  \label{fig:extruder}
\end{figure}

\section{Tier 2: perturbations through the bounce, and the CMB}
\label{sec:tier2}

This is where the framework meets its central observable, and a first pass already
clears the main hurdle. Evolving tensor modes through the bounce with the
Mukhanov--Sasaki equation~\citep{mukhanov1992}, with Bunch--Davies initial data
deep in contraction,
reproduces the analytic matter-bounce result~\citep{wands1999}: for a contraction
of equation of
state $w$ the tensor tilt is $n_T=3-|2p-1|$ with $p=2/(1+3w)$, and a
matter-dominated contraction ($w=0$) gives $n_T=0$ \emph{exactly}---a
scale-invariant primordial spectrum, the observed near-scale-invariance,
\textbf{with no inflaton} (\Cref{fig:primordial}). What Einstein--Cartan supplies
is the one ingredient the matter-bounce scenario otherwise
lacks~\citep{novello2008,brandenberger2017}: a genuinely
non-singular bounce, which rival realisations buy with exotic ghost matter but
torsion delivers for free. The slight observed red tilt $n_s\simeq0.965$
corresponds to a contraction just softer than matter ($w\approx-0.003$).

The acoustic peaks then follow from \emph{shared} physics: the framework runs the
entire post-bounce hot Big Bang exactly as $\Lambda$CDM does---same plasma, same
recombination at $z\simeq1090$, same sound horizon---so it inherits the peaks
rather than re-deriving them. Evaluated on the standard background the comoving
sound horizon is $r_s\simeq142\,\mathrm{Mpc}$ and the peak series
$\ell_n\simeq230,540,850,1150,1460$ matches the observed Planck peaks
($220,540,810,1120,1450$) to a few percent (\Cref{fig:acoustic}). The hard,
decisive number---the acoustic scale---is reproduced; what remains is a full
Boltzmann computation (CAMB/CLASS-class) of peak \emph{heights}, sensitive to the
baryon and dark-matter densities and the precise primordial amplitude. The
parent-Hawking reading of the CMB is consistent with all of this: it concerns the
\emph{origin} of the thermal plasma at the bounce surface, not a replacement for
recombination---the peaks arise from standard acoustics either way.

\begin{figure}[t]
  \centering
  \includegraphics[width=\textwidth]{primordial}
  \caption{The primordial spectrum from the Einstein--Cartan bounce. \emph{Left:}
    the tensor power spectrum is flat (scale-invariant) for a matter-dominated
    contraction and blue for radiation---scale invariance with no inflaton.
    \emph{Right:} the tilt $n_T$ versus contraction equation of state on the
    analytic curve; $w=0$ is scale-invariant, and the observed tilt corresponds to
    $w\approx-0.003$. Reproducible as \texttt{figures/scripts/ch09\_primordial.py}.}
  \label{fig:primordial}
\end{figure}

\begin{figure}[t]
  \centering
  \includegraphics[width=\textwidth]{acoustic}
  \caption{The acoustic peaks. The bounce's scale-invariant primordial input,
    carried through the standard post-bounce plasma, yields a sound horizon
    $r_s\simeq142\,\mathrm{Mpc}$ and a peak series matching the observed Planck
    multipoles to a few percent. The peak \emph{heights} await a full Boltzmann
    computation. Reproducible as \texttt{figures/scripts/ch09\_acoustic.py}.}
  \label{fig:acoustic}
\end{figure}

\section{Tiers 3--5, and the frontier}

The higher tiers are multi-year programs, sketched here as scope rather than
results. \textbf{Tier~3 (chirality dynamics)} would quantify the membrane's filter
efficiency and predict the dark-to-baryon ratio's $f\approx1/6$ from
Einstein--Cartan microphysics---turning a consistency target into a derivation.
\textbf{Tier~4 (population statistics)} treats the supraverse as a stochastic
process---nucleation, growth, death, recursion---and computes the equilibrium
distribution whose typical observer Chapter~\ref{chap:nowobs} located, supplying
the rigorous measure the Boltzmann-brain argument (Chapter~\ref{chap:void}) still
owes. \textbf{Tier~5} feeds all of this into direct, quantitative comparisons with
specific datasets. Further out sit the genuinely frontier tests: gravitational
decoherence---the Di\'osi--Penrose collapse
proposal~\citep{diosi1989,penrose1996}, among the testable collapse
models~\citep{bassi2013}---as the laboratory-scale limb of the same
gravity-separates-branches
physics as the bounce, information-theoretic signatures of geometric dark matter,
and any direct gravitational imprint of the parent. These are catalogued for
completeness; none is near-term.

\section{The minimum viable project, and the ethos}

If one project should come first, it is \textbf{Tier~2 with a focus on the CMB
power spectrum}: decisive either way, publishable either way, bounded at about a
year, and within reach of standard numerical-relativity-plus-perturbation-theory
expertise. The visualization track (Chapter~\ref{chap:almostnothing}) can proceed
in parallel, needing no Tier-2 output to begin.

The governing principle is not to ``prove'' the framework by simulation but to
make it \emph{testable}: code open from the start, validated against analytic
limits, checked for numerical convergence, and reproducible by anyone. That
standard is already met by the results above---every figure and number in this
book regenerates from the repository, and the simulation library ships with a
passing test suite. Reproducible code lets others check, extend, and challenge the
claims, which is the only thing that turns a worldview into science.
